\par\Needspace{6\baselineskip}\medskip
\hypertarget{native-076-article058}{}\pdfbookmark[2]{58.}{article-058}
\begin{center}58.\end{center}
Unsere Formeln für die Abstände von gegebenen Ebenen enthalten $v$ und $r$. Wenn man diese Größen erst aus der Zeit bestimmen muss, kann man einen Teil der Rechenschritte noch abkürzen und den Aufwand dadurch merklich verringern. Die Abstände lassen sich nämlich mit einer sehr einfachen Formel unmittelbar aus der exzentrischen Anomalie in der Ellipse oder aus der Hilfsgröße $F$ oder $u$ in der Hyperbel bestimmen. Damit entfällt die Berechnung der wahren Anomalie und des Radiusvektors ganz. Der Ausdruck $kr\sin(v+K)$ wird wie folgt umgeformt.

I. \textit{Für die Ellipse} wird er unter Beibehaltung der Bezeichnungen aus Artikel 8 zu
\[
ak\cos\varphi\cos K\sin E+ak\sin K(\cos E-e).
\]
Bestimmt man also $l$, $L$, $\lambda$ durch die Gleichungen
\[
\begin{aligned}
ak\sin K&=l\sin L\\
ak\cos\varphi\cos K&=l\cos L\\
-eak\sin K&=-el\sin L=\lambda,
\end{aligned}
\]
\SourceLeaf{077}so geht unser Ausdruck in $l\sin(E+L)+\lambda$ über. Dabei sind $l$, $L$, $\lambda$ konstant, soweit man $k$, $K$, $e$ als konstant ansehen kann. Andernfalls gilt für die Berechnung ihrer Änderungen dasselbe wie im vorigen Artikel.

Als Beispiel geben wir die Umformung des in Artikel 56 gefundenen Ausdrucks für $x$ an. Wir setzen die Länge des Perihels $=121^0\,17^{\prime}\,34^{\prime\prime}{,}4$, $\varphi=14^0\,13^{\prime}\,31^{\prime\prime}{,}97$ und $\log a=0{,}442\,3790$. Damit ergibt sich der Abstand des Perihels vom aufsteigenden Knoten in der Ekliptik zu $=308^0\,49^{\prime}\,20^{\prime\prime}{,}7=u-v$. Daraus folgt $K=212^0\,36^{\prime}\,56^{\prime\prime}{,}09$. Somit erhalten wir
\begin{center}
\begin{tabular}{@{}ll@{}}
\makebox[13em][l]{$\log ak$\hfill .\,.\,.\,.\,.\,.\,.} & $0{,}441\,1713$\\[.3em]
\makebox[13em][l]{$\log\sin K$\hfill .\,.\,.\,.\,.} & $9{,}731\,5887_n$\\[.3em]
\makebox[13em][l]{$\log ak\cos\varphi$\hfill .\,.\,.} & $0{,}427\,6456$\\[.3em]
\makebox[13em][l]{$\log\cos K$\hfill .\,.\,.\,.\,.} & $9{,}925\,4698_n$
\end{tabular}
\end{center}
\begin{center}
\begin{tabular}{@{}ll@{}}
\makebox[13em][l]{$\log l\sin L$\hfill .\,.\,.\,.} & $0{,}172\,7600_n$\\[.3em]
\makebox[13em][l]{$\log l\cos L$\hfill .\,.\,.\,.} & $0{,}353\,1154_n$\\[.3em]
\hline
daraus $L=$ & $213^0\,25^{\prime}\,51^{\prime\prime}{,}31$\\[.3em]
\makebox[13em][l]{$\log l$\hfill .\,.\,.\,.\,.\,.\,.\,.} & $0{,}431\,6627$\\[.3em]
\makebox[13em][l]{$\log\lambda$\hfill .\,.\,.\,.\,.\,.\,.\,.} & $9{,}563\,2352$\\[.3em]
$\lambda=$ & $+0{,}365\,7929$
\end{tabular}
\end{center}

II. Für die Hyperbel wird die Formel $kr\sin(v+K)$ nach Artikel 21 zu
\[
\lambda+\mu\operatorname{tang}F+\nu\operatorname{secans}F,
\]
wenn man setzt
\[
ebk\sin K=\lambda,\qquad bk\operatorname{tang}\psi\cos K=\mu,\qquad -bk\sin K=\nu;
\]
Offenbar lässt sich derselbe Ausdruck auch auf die Form $\frac{n\sin(F+N)+\nu}{\cos F}$ bringen. Verwendet man anstelle von $F$ die Hilfsgröße $u$, so geht der Ausdruck $kr\sin(v+K)$ nach Artikel 21 in $\alpha+\beta u+\frac{\gamma}{u}$ über. Dabei werden $\alpha$, $\beta$, $\gamma$ durch die Formeln bestimmt
\[
\begin{aligned}
\alpha&=\lambda=ebk\sin K\\
\beta&=\tfrac{1}{2}(\nu+\mu)=-\tfrac{1}{2}ebk\sin(K-\psi)\\
\gamma&=\tfrac{1}{2}(\nu-\mu)=-\tfrac{1}{2}ebk\sin(K+\psi).
\end{aligned}
\]

III. In der Parabel, wo sich die wahre Anomalie unmittelbar aus der Zeit ergibt, braucht man nur noch den Wert des Radiusvektors einzusetzen. Bezeichnet man also den Abstand im Perihel mit $q$, so wird der Ausdruck $kr\sin(v+K)$ zu $=\frac{qk\sin(v+K)}{\cos\tfrac{1}{2}v^2}\mathord{\cdot}$
